\documentclass[conference]{IEEEtran}
\IEEEoverridecommandlockouts
% The preceding line is only needed to identify funding in the first footnote.
\usepackage{cite}
\usepackage{amsmath,amssymb,amsfonts}
\usepackage{algorithmic}
\usepackage{graphicx}
\usepackage{textcomp}
\usepackage{xcolor}
\def\BibTeX{{\rm B\kern-.05em{\sc i\kern-.025em b}\kern-.08em
    T\kern-.1667em\lower.7ex\hbox{E}\kern-.125emX}}

% --- Helper for visible placeholders / TODOs (remove before submission) ---
\newcommand{\TODO}[1]{\textcolor{red}{[\textbf{TODO:} #1]}}
\newcommand{\PLACEHOLDER}[1]{\textcolor{red}{\textit{[Placeholder --- #1]}}}

\begin{document}

\title{Risk-Aware Distributional Reinforcement Learning\\for 5G LTM Handover Decisions}

\author{\IEEEauthorblockN{1\textsuperscript{st} Given Name Surname}
\IEEEauthorblockA{\textit{dept. name of organization (of Aff.)} \\
\textit{name of organization (of Aff.)}\\
City, Country \\
email address or ORCID}
\and
\IEEEauthorblockN{2\textsuperscript{nd} Given Name Surname}
\IEEEauthorblockA{\textit{dept. name of organization (of Aff.)} \\
\textit{name of organization (of Aff.)}\\
City, Country \\
email address or ORCID}
\and
\IEEEauthorblockN{3\textsuperscript{rd} Given Name Surname}
\IEEEauthorblockA{\textit{dept. name of organization (of Aff.)} \\
\textit{name of organization (of Aff.)}\\
City, Country \\
email address or ORCID}
}

\maketitle

% =====================================================================
\begin{abstract}
Handover (HO) decisions in 5G lower-layer-triggered mobility (LTM) must
balance throughput against reliability under fast-varying radio conditions,
where rare but costly events---handover and radio-link failures---dominate
the user experience. We cast LTM HO as a Markov decision process and study
distributional reinforcement learning (QR-DQN) with risk-aware action
selection. We report two separable and statistically significant gains over
scalar DQN: a large gain from the distributional learning target itself, and
a second gain from a Conditional-Value-at-Risk (CVaR) policy that optimises
the lower tail of the return distribution. A \emph{de-confounded} sweep that
holds the number of predicted quantiles fixed while varying the risk level
reveals an interior risk optimum at $\mathrm{CVaR}$ level $\alpha\!\approx\!0.5$,
reversing the conclusion of a naive sweep in which capacity and risk are
entangled. On a 1{,}000-trajectory 5G LTM benchmark the risk-aware agent
nearly halves the handover-failure rate of DQN, surpasses a hand-tuned LTM
baseline---evaluated in the same simulator---on capacity, reliability and
ping-pong rate, and improves on a published contextual-bandit baseline in
capacity and reliability. A quantile-bridge
ablation shows both ingredients are necessary: a single risk-neutral quantile
collapses to DQN, and a single risk-averse quantile only reaches risk-neutral
performance---the full gain requires a multi-quantile CVaR representation.
\end{abstract}

\begin{IEEEkeywords}
distributional reinforcement learning, quantile regression, conditional value
at risk, risk-sensitive control, 5G NR, lower-layer triggered mobility,
handover optimization.
\end{IEEEkeywords}

% =====================================================================
\section{Introduction}

Fifth-generation (5G) networks densify the radio access with many small,
overlapping cells, so a moving user equipment (UE) crosses cell boundaries
frequently and must be handed over from one cell---or sector---to another many
times during a session. Each handover (HO) is a brief, risk-laden
reconfiguration: chosen well it preserves throughput and reliability; chosen
poorly it triggers a handover failure (HOF), a radio-link failure (RLF), or an
unnecessary ``ping-pong'' (PP) back to the previous cell. \emph{Lower-layer
triggered mobility} (LTM) is the 5G New Radio mechanism that accelerates these
decisions by triggering them from fast L1/L2 measurements rather than slower
higher-layer signalling; the price of the extra speed is that the choice of
\emph{which} cell/sector to connect to at each measurement opportunity must be
made under rapidly varying radio conditions. That selection is the decision
problem we study.

Learning-based HO optimisation typically casts cell selection as a heuristic
rule or a contextual-bandit / scalar value-learning problem that maximises an
\emph{expected} objective---throughput, or a weighted KPI score---sometimes
augmented with reconfigurable intelligent surfaces (RIS) and LMMSE channel
prediction to improve the underlying signal quality~\cite{ainna}. Expected-value
objectives are risk-neutral: they optimise the \emph{average} outcome and are
indifferent to its dispersion. Yet in LTM the user experience is dominated not
by the average step but by rare, disproportionately costly tail
events---handover and radio-link failures---which a mean-optimal policy has no
explicit incentive to avoid. This motivates a \emph{risk-aware} formulation
that acts on the lower tail of the outcome distribution rather than its mean.

In this paper we approach LTM HO with \emph{distributional} reinforcement
learning, which models the full return distribution $Z(x,a)$ rather than its
mean $Q(x,a)$, and with \emph{risk-aware} action selection based on the
Conditional Value at Risk (CVaR) of that distribution. Intuitively, learning
the distribution yields a richer training signal, and acting on its lower tail
makes the agent avoid worst-case outcomes (failures) instead of merely
maximising the average. Our contributions are:

\begin{itemize}
  \item We formulate LTM HO as a Markov decision process and apply
  Quantile-Regression DQN (QR-DQN) on a calibrated 5G LTM simulator
  (Section~\ref{sec:model}--\ref{sec:method}).
  \item We introduce a \emph{truncated-CVaR} head that concentrates network
  capacity on only the lower $\alpha$-tail of the return distribution that
  actually drives a CVaR policy (Section~\ref{sec:method}).
  \item We \emph{de-confound} the two coupled knobs of the risk-aware
  agent---the number of predicted quantiles $k$ (capacity) and the CVaR level
  $\alpha$ (risk)---and show an interior risk optimum at $\alpha^\star\!\approx\!0.5$
  that a naive, capacity-entangled sweep hides (Section~\ref{sec:results-sweep}).
  \item We isolate where the gain originates with a \emph{quantile-bridge}
  ablation, demonstrating that both the multi-quantile representation and the
  CVaR policy are individually necessary (Section~\ref{sec:results-bridge}).
  \item On the 5G LTM benchmark we demonstrate two separable, significant gains
  over DQN, and benchmark the risk-aware policy against both a hand-tuned LTM
  baseline (run in our own simulator) and the contextual multi-armed bandit
  (CMAB) of~\cite{ainna}.
\end{itemize}

% =====================================================================
\section{Background and Related Work}
\label{sec:related}

\subsection{LTM and Handover Optimization}
In legacy mobility the network reconfigures a UE's serving cell through
higher-layer (RRC) signalling driven by filtered L3 measurements---robust but
slow. \emph{Lower-layer triggered mobility} (LTM) instead prepares one or more
candidate cells in advance and lets a fast L1/L2 measurement event trigger the
execution of a pre-configured switch, cutting the interruption time but
tightening the decision loop. The operational cost of a mobility policy is
captured by a small set of events: a \emph{handover} (HO) whenever the serving
cell changes, a \emph{ping-pong} (PP) when the UE returns to a recently left
cell within a short window, a \emph{handover failure} (HOF) when the target
link cannot be established, and a \emph{radio-link failure} (RLF) when the
serving link collapses after a sustained out-of-sync period. A good policy
minimises failures and ping-pongs while sustaining spectral efficiency. Prior
learning-based work optimises this trade-off with heuristics and contextual
bandits over an expected-KPI objective, optionally stacking RIS and LMMSE
channel prediction on top to improve signal quality and target
pre-selection~\cite{ainna}; this is the closest prior art to ours. It is,
however, risk-neutral---the gap we address.

\subsection{Distributional Reinforcement Learning}
\label{sec:rel-distrl}
Distributional RL predicts the full distribution of the return rather than its
expectation. Instead of the scalar action-value $Q(x,a)$, the agent models the
return as a random variable $Z(x,a)$ that obeys the distributional Bellman
equation
\begin{equation}
Z(x,a) \overset{D}{=} R(x,a) + \gamma\,Z(X',A'),
\label{eq:dist-bellman}
\end{equation}
where $\overset{D}{=}$ denotes equality in distribution, $R(x,a)$ is the
reward, $\gamma$ the discount factor, and $(X',A')$ the (random) next
state--action pair. Even when the policy ultimately acts on the mean
$\mathbb{E}[Z(x,a)]$, learning the distribution provides a richer per-update
signal and consistently improves the resulting policy. Most practical
algorithms approximate the cumulative distribution as a staircase function;
representative methods include C51~\cite{c51}, QR-DQN~\cite{qrdqn},
IQN~\cite{iqn} and FQF~\cite{fqf}. We build on QR-DQN, which parametrises the
distribution by the locations of a fixed set of quantiles.

\subsection{Risk-Sensitive Reinforcement Learning}
\label{sec:rel-risk}
In standard distributional RL the greedy action maximises
$\mathbb{E}[Z(x,a)]$, approximated as the average of the $N$ predicted quantiles
$\tfrac{1}{N}\sum_{i=1}^{N}\theta_i(x,a)$. One may instead aggregate the
$\theta_i$ under a different criterion to obtain a risk-sensitive policy; note
that no aggregation can totally order all distributions~\cite{rockafellar}.
A common risk-averse choice replaces the expectation with the Conditional
Value at Risk at level $\alpha$, $\mathrm{CVaR}_\alpha$, the average of the
worst $\alpha$-fraction of outcomes. For a uniform return on $[0,1]$, for
example, $\mathrm{CVaR}_{0.2}$ is the mean of the values in $[0,0.2]$, i.e.\
$0.1$. In problems such as LTM HO, risk-aware decision-making is attractive
because avoiding worst-case outcomes (handover or radio-link failures) is more
valuable than maximising the expected return alone. CVaR and related
risk-sensitive criteria have a long history in reinforcement learning and
operations research~\cite{chow2014,tamar2015}; distributional RL makes them
particularly convenient, since the return quantiles needed to evaluate
$\mathrm{CVaR}_\alpha$ are predicted directly by the network.

% =====================================================================
\section{System Model and Problem Formulation}
\label{sec:model}

\subsection{5G LTM Deployment and Simulator}
We evaluate on a calibrated multi-sector 5G deployment of $7$ tri-sector base
stations ($21$ sectors), driven by a dataset of $1{,}000$ pre-computed UE
trajectories of $300$\,s each whose per-sector channel gains are sampled every
$10$\,ms. The simulator advances at this $10$\,ms resolution while the agent
issues one HO decision every $100$\,ms. At each step the serving
signal-to-interference-plus-noise ratio (SINR) is computed from the channel
gains, transmit power and inter-cell interference, then mapped to a
modulation-and-coding scheme (MCS) and hence a spectral efficiency through a
$26$-entry SINR$\to$MCS table whose lowest entry (SINR $<-3$\,dB) is an outage
(MCS$=0$); a sustained run of outage steps escalates to a radio-link failure.
The radio and mobility constants (Table~\ref{tab:sysparams}) are calibrated to
the reference LTM implementation of~\cite{ainna} rather than to nominal
standards values, so that our LTM baseline reproduces that work's published
KPIs; the same physics module is shared by the RL environment and a standalone
reference simulator used for parity checking.

\begin{table}[htbp]
\caption{Key simulator constants (calibrated to the reference LTM implementation).}
\label{tab:sysparams}
\begin{center}
\begin{tabular}{|l|c|}
\hline
\textbf{Parameter} & \textbf{Value} \\
\hline
Base stations $\times$ sectors      & $7 \times 3 = 21$ \\
UE trajectories                     & $1{,}000$ ($300$\,s each) \\
Simulator / decision step           & $10$\,ms / $100$\,ms \\
Transmit power                      & $25$\,dBm \\
Noise floor                         & $-101$\,dBm ($-174$\,dBm/Hz $\times\,20$\,MHz) \\
Receiver sensitivity                & $-95$\,dBm \\
Outage (RLF) SINR threshold         & $-3$\,dB \\
Execution power offset              & $3.0$\,dB \\
Max.\ prepared cells                & $5$ \\
\hline
\end{tabular}
\end{center}
\end{table}

\subsection{MDP Formulation}
\label{sec:mdp}
We model LTM HO as a Markov decision process.

\textbf{State.} An $88$-dimensional vector: UE speed $(1)$ and serving-cell
tenure $(1)$, the one-hot serving sector $(21)$, per-sector RSRP $(21)$,
moving-average MCS $(21)$ and moving-average SINR $(21)$, and the UE position
$(x,y)$ $(2)$.

\textbf{Action.} At each $100$\,ms step the agent selects which of the $21$
sectors to connect to. If the chosen sector equals the current serving sector
no handover is performed; otherwise a handover is triggered.

\textbf{Reward.} The agent optimises the multiplicative LTM HO reward
of~\cite{ainna},
\begin{equation}
R_t = R^{\mathrm{thr}}_t \,\cdot\,
      \alpha_{\mathrm{HO}}^{\,\mathbb{1}^{\mathrm{HO}}_t}\,\cdot\,
      \alpha_{\mathrm{PP}}^{\,\mathbb{1}^{\mathrm{PP}}_t}\,\cdot\,
      \alpha_{\mathrm{HOF}}^{\,\mathbb{1}^{\mathrm{HOF}}_t}\,\cdot\,
      \rho_t,
\label{eq:reward}
\end{equation}
where $R^{\mathrm{thr}}_t$ is the throughput term (the serving spectral
efficiency in bps/Hz); the indicators
$\mathbb{1}^{\mathrm{HO}}_t,\mathbb{1}^{\mathrm{PP}}_t,\mathbb{1}^{\mathrm{HOF}}_t\in\{0,1\}$
flag a handover, ping-pong and handover failure at step $t$; and
$\alpha_{\mathrm{HO}},\alpha_{\mathrm{PP}},\alpha_{\mathrm{HOF}}\in(0,1]$ are
the corresponding multiplicative penalties (defaults $\alpha_{\mathrm{HOF}}{=}0.1$,
$\alpha_{\mathrm{HO}}{=}0.8$, $\alpha_{\mathrm{PP}}{=}0.9$, so a failure is
penalised far more harshly than a routine handover). The reliability factor
$\rho_t\in(0,1]$ is a soft reverse-sigmoid of the out-of-sync counter that
decays the reward as the serving link approaches radio-link failure. Since
every factor lies in $(0,1]$, the reward is the fraction of attainable
throughput that the policy retains after instability, which the distributional
agent models as a random return.

% =====================================================================
\section{Method}
\label{sec:method}

\subsection{QR-DQN and the Quantile Loss}
QR-DQN parametrises $Z(x,a)$ by its inverse CDF at $N$ quantile fractions
$\tau_1,\dots,\tau_N\in[0,1]$, learning
$\theta_i(x,a)\approx F^{-1}_{Z(x,a)}(\tau_i)$. The aggregated value used for
action selection,
\begin{equation}
\mathbb{E}[Z(x,a)] = \int_0^1 F^{-1}_{Z(x,a)}(\tau)\,d\tau,
\label{eq:expectation}
\end{equation}
is replaced by a finite-sample quadrature $\sum_{i=1}^{N} w_i\,\theta_i(x,a)$
with weights $w_i\ge 0$, $\sum_i w_i = 1$. The choice of $(\tau_i,w_i)$ defines
the quantile-positioning scheme. The quantiles are trained with the quantile
(pinball) Huber loss. Writing the Huber function with threshold $\kappa$ as
\begin{equation}
\mathcal{H}_\kappa(u) =
\begin{cases}
\tfrac12 u^2, & |u|\le\kappa,\\[2pt]
\kappa\big(|u|-\tfrac12\kappa\big), & |u|>\kappa,
\end{cases}
\end{equation}
the quantile-Huber penalty at fraction $\tau$ is
$\rho^\kappa_\tau(u) = \big|\tau - \mathbb{1}\{u<0\}\big|\,
\mathcal{H}_\kappa(u)/\kappa$, and the per-transition loss sums the pairwise
temporal-difference errors between the $N$ predicted quantiles and the $N$
Bellman-target atoms $y_j = r + \gamma\,\theta_j\!\big(x',a^\star\big)$,
\begin{equation}
\mathcal{L}(x,a) = \sum_{i=1}^{N}\sum_{j=1}^{N}
   \tilde w_j\,\rho^{\kappa}_{\tau_i}\!\big(y_j - \theta_i(x,a)\big),
\label{eq:qr-loss}
\end{equation}
with $a^\star$ the greedy action under the aggregated value and $\tilde w_j$
the target-axis quadrature weights ($\tilde w_j{=}1/N$ for the midpoint rule).
At $N{=}1,\,\tau{=}0.5$ the asymmetric weight $|\tau-\mathbb{1}\{u<0\}|$ becomes
the constant $\tfrac12$ and \eqref{eq:qr-loss} reduces to a symmetric Huber
regression onto the median---i.e.\ (Huberised) scalar value learning---a fact
we exploit in Section~\ref{sec:results-bridge}. We use $\kappa{=}0.25$
throughout.

\subsection{Quantile-Positioning Schemes}
QR-DQN papers default to the midpoint rule; we additionally consider two
classical quadrature alternatives and ask whether either the integration
accuracy of the aggregator or the placement of the learnable quantiles affects
policy quality on LTM HO.

\textbf{Midpoint rule.} Vanilla QR-DQN places $\tau_i=(i-\tfrac12)/N$ with
uniform weights $w_i=1/N$---the midpoint Riemann sum for
$\int_0^1 F^{-1}(\tau)\,d\tau$, with $O(1/N^2)$ error for smooth integrands.
All final models use this scheme; Section~\ref{sec:results-qmode} confirms the
two alternatives below give no significant gain on LTM.

\textbf{Gauss--Legendre quadrature.} Nodes are placed at the roots of the
$N$-th Legendre polynomial shifted to $[0,1]$, integrating polynomials up to
degree $2N-1$ exactly. The network still outputs $N$ quantiles, but the
$\tau$ used in the loss and the aggregation weights follow the Gauss--Legendre
rule, so the parameter count is unchanged.

\textbf{Trapezoidal rule with fixed endpoints.} Uniform spacing including the
endpoints $\tau_0=0,\tau_{N-1}=1$, with $w_0=w_{N-1}=1/(2(N-1))$ and
$w_i=1/(N-1)$ otherwise. Because $F^{-1}$ is unbounded at the endpoints for
infinite-support distributions, we fix them a priori from physical
considerations ($F^{-1}(0)=0$ since returns are non-negative under the
multiplicative reward, and $F^{-1}(1)=q_{\max}$ a conservative upper bound),
predicting only the $N-2$ interior quantiles.

\subsection{Risk-Aware Selection and Truncated CVaR}
\label{sec:trunc-cvar}
A $\mathrm{CVaR}_\alpha$ policy integrates only over $\tau\le\alpha$: the upper
$1-\alpha$ of a predicted quantile vector is generated and then discarded---wasted
capacity. We instead train the network on only $k=\lceil\alpha N\rceil$
quantiles placed uniformly in $[0,\alpha]$, removing the upper lobe from both
the prediction and the loss. Action selection is then the mean of the $k$
predictions, which under the midpoint rule is exactly $\mathrm{CVaR}_\alpha$.
Concentrating capacity on the part of the distribution that drives the policy
should yield a better CVaR estimate at lower cost; Section~\ref{sec:results-trunc}
confirms this directly---the full-grid alternative, given the identical policy,
is both significantly worse and slower to train.

% =====================================================================
\section{Experimental Setup}
\label{sec:setup}

\subsection{Evaluation Protocol and Metrics}
All agents are trained and evaluated on the same $1{,}000$ pre-computed UE
trajectories ($300$\,s each, $10$\,ms simulation resolution, one HO decision
every $100$\,ms). We deliberately use the full trajectory set for both
training and final evaluation rather than a held-out split: the channel-gain
dataset is fixed and the policy never observes future returns, so the relevant
question is how well each agent exploits the available propagation environment,
not generalisation to unseen geometry. Final numbers freeze the learned
weights, set $\epsilon=0$ (greedy), and replay all $1{,}000$ UEs. Every
configuration is run for $5$ random seeds $\times$ $2{,}000$ episodes; we
report the cross-seed mean $\pm$ standard deviation. Alongside the reward we
track the eight physical KPIs of~\cite{ainna}: handover rate (HO/min),
handover-failure rate (HOF/min), ping-pong rate (PP/min), radio-link-failure
rate (RLF/min), spectral efficiency / capacity (bps/Hz), link reliability
(\%), cell preparations (/min) and resource reservation (\%).

\subsection{Models and Hyperparameters}
We compare three learned agents that share an identical two-layer MLP trunk
($256\times256$) and the same observation, action and reward definitions, so
that any difference is attributable to the \emph{learning target} rather than
capacity or input:
\begin{enumerate}
  \item \textbf{DQN}---scalar value learning; greedy $\arg\max_a Q(x,a)$.
  \item \textbf{QR-DQN (risk-neutral, RN)}---distributional with a mean policy
  $\mathbb{E}[Z(x,a)]$; predicts $10$ quantiles over $[0,1]$.
  \item \textbf{QR-DQN (risk-aware, RA)}---distributional with a
  $\mathrm{CVaR}_{0.5}$ policy via the truncated scheme of
  Section~\ref{sec:trunc-cvar}: it predicts only the $k=\lceil\alpha N\rceil=10$
  lower quantiles ($\alpha=0.5$, $N=20$).
\end{enumerate}
The two QR-DQN variants are matched to the \emph{same} effective number of
predicted quantiles ($k=10$), hence the same head size and parameter count;
they differ only in \emph{where} the quantiles sit (full $[0,1]$ vs.\ lower
half $[0,0.5]$) and the selection rule (mean vs.\ CVaR). All three share a
single soft-target-update backbone ($\tau=0.01$); the hard-vs-soft choice is
within cross-seed noise for both agents, so a single shared scheme keeps the
comparison strictly apples-to-apples. Hyperparameters are in
Table~\ref{tab:hparams}; the learning rate, batch size, discount and
target-update rate were fixed by coarse grid searches and partly inherited from
the prior work (halving the learning rate to $5{\times}10^{-5}$ is $2.3\sigma$
worse at the $2000$ ep $\times$ $5$ seed budget), and preliminary experiments
with $n$-step returns, gradient
clipping and optimiser-$\epsilon$ variants were not adopted as they did not
improve the LTM agents (on the final risk-aware agent, $3$-step returns are
$3.4\sigma$ worse and markedly less stable). Two choices the prior work leaves
implicit we verified directly on the final backbone ($1000$ ep $\times$ $3$
seeds). First, the Huber threshold is immaterial: across
$\kappa\in\{0.1,0.25,0.5,1.0\}$ reward spans only $1.5$ points
($1081.1$--$1082.6$), within cross-seed noise (Table~\ref{tab:kappa}), so we
keep the standard $\kappa{=}0.25$. Second,
network capacity saturates at $256{\times}256$: halving the width to
$128{\times}128$ costs $5.4\sigma$ of reward ($1076.6$ vs.\ $1082.1$), while
doubling it to $512{\times}512$ adds nothing ($1082.3$; $+0.3\sigma$).

\begin{table}[htbp]
\caption{Agent hyperparameters (shared backbone).}
\label{tab:hparams}
\begin{center}
\begin{tabular}{|l|c|c|c|}
\hline
\textbf{Hyperparameter} & \textbf{DQN} & \textbf{QR-DQN (RN)} & \textbf{QR-DQN (RA)} \\
\hline
Hidden layers      & $256\times256$ & $256\times256$ & $256\times256$ \\
Learning rate      & $10^{-4}$ & $10^{-4}$ & $10^{-4}$ \\
Discount $\gamma$  & 0.9 & 0.9 & 0.9 \\
Replay buffer      & $10^5$ & $10^5$ & $10^5$ \\
Batch size         & 256 & 256 & 256 \\
Train frequency    & 1 & 1 & 1 \\
Target update      & soft $\tau{=}0.01$ & soft $\tau{=}0.01$ & soft $\tau{=}0.01$ \\
$\epsilon$ sched.  & $1.0\!\to\!0.05$ & idem & idem \\
Huber $\kappa$     & --- & 0.25 & 0.25 \\
Nominal $N$        & --- & 10 & 20 \\
Predicted $k$      & --- & 10 & 10 \\
Risk policy        & --- & mean & $\mathrm{CVaR}_{0.5}$ \\
Quantile support   & --- & $[0,1]$ & $[0,0.5]$ \\
\hline
\end{tabular}
\end{center}
\end{table}

\begin{table}[htbp]
\caption{Huber threshold $\kappa$ on the final backbone (truncated
$\mathrm{CVaR}_{0.5}$, $k{=}10$; $1000$ ep $\times$ $3$ seeds; mean $\pm$ std on
reward). Reward spans only $1.5$ points across an order of magnitude in
$\kappa$ (within cross-seed noise); we keep the standard $\kappa{=}0.25$ (bold).}
\label{tab:kappa}
\begin{center}
\setlength{\tabcolsep}{4pt}
\begin{tabular}{|l|c|c|c|c|}
\hline
$\kappa$ & $0.10$ & $\mathbf{0.25}$ & $0.50$ & $1.00$ \\
\hline
Reward $\uparrow$    & $1081.3{\pm}0.6$ & $\mathbf{1082.1{\pm}1.0}$ & $1082.6{\pm}1.9$ & $1081.1{\pm}1.7$ \\
HOF/min $\downarrow$ & $1.36$ & $\mathbf{1.28}$ & $1.31$ & $1.35$ \\
\hline
\end{tabular}
\end{center}
\end{table}

% =====================================================================
\section{Results on the LTM Handover Task}
\label{sec:results}

\subsection{Choosing the Risk Geometry: a De-confounded Sweep}
\label{sec:results-sweep}
The risk-aware agent couples two knobs: the number of in-use quantiles $k$ and
the CVaR level $\alpha$. Because $k=\lceil\alpha N\rceil$, a naive sweep that
fixes the \emph{nominal} $N$ and varies $\alpha$ silently changes $k$ too,
conflating ``how risk-averse'' with ``how many quantiles.'' We therefore sweep
the two axes separately ($1{,}000$ episodes $\times$ $3$ seeds per point).

\textbf{Quantile count.} Fixing $\alpha=0.5$ and varying the head so
$k\in\{10,13,15,25,50\}$ leaves reward essentially flat---$1083.8$, $1084.6$,
$1083.4$, $1082.4$, $1083.1$ (all within cross-seed noise)---so we adopt the
small, cheap $k=10$ head.

\textbf{Risk level.} Holding $k=10$ and varying $\alpha$ (adjusting $N$ so
$\lceil\alpha N\rceil=10$ at every point) traces a concave curve with an
interior optimum (Table~\ref{tab:alpha}). Reward peaks at $\alpha\approx0.5$
atop a broad plateau ($\alpha\in[0.3,0.7]$ all within $\sim\!1081$--$1084$) and
falls sharply toward the risk-neutral extreme: $\alpha=1.0$ (which, since
$\mathrm{CVaR}_{1.0}=\mathbb{E}[Z]$, \emph{is} the risk-neutral policy) is
$9.2$ reward points---about $7\sigma$---below $\alpha=0.5$, and simultaneously
worst on every safety KPI (HOF $1.83$ vs.\ $1.28$, RLF $0.148$ vs.\ $0.103$).
This reverses what a confounded $N$-fixed sweep suggests: once the quantile
count is held fixed, a genuine risk-aware optimum near $\alpha\approx0.5$
emerges. We adopt $\alpha^\star=0.5$---CVaR of the worse half.

\begin{table}[htbp]
\caption{Reward vs.\ CVaR level $\alpha$ at fixed $k=10$ ($1000$ ep $\times$ $3$ seeds).}
\label{tab:alpha}
\begin{center}
\begin{tabular}{|l|c|c|c|c|c|}
\hline
$\alpha$ & 0.30 & 0.40 & 0.50 & 0.70 & 1.00 \\
\hline
Reward & $1081.2$ & $1081.5$ & $\mathbf{1083.8}$ & $1081.2$ & $1074.6$ \\
\hline
\end{tabular}
\end{center}
\end{table}

\subsection{Quantile Positioning Does Not Matter on LTM}
\label{sec:results-qmode}
Section~\ref{sec:method} asked whether the placement of the learnable quantiles
or the accuracy of the aggregating quadrature affects policy quality. We test
this directly, training the risk-neutral agent ($N{=}10$, mean policy) under the
midpoint, Gauss--Legendre and trapezoidal schemes with everything else held
fixed ($1000$ ep $\times$ $3$ seeds, final backbone). The three are
\emph{statistically indistinguishable} (Table~\ref{tab:qmode}): relative to
midpoint, Gauss--Legendre shifts reward by $+0.8$ ($0.5\sigma$) and trapezoidal
by $+2.8$ ($1.3\sigma$), both within cross-seed noise, and the safety KPIs
differ by at most ${\approx}2\sigma$. Neither the higher-order accuracy of
Gauss--Legendre nor the variance reduction of the fixed trapezoidal
endpoints\footnote{The trapezoidal upper endpoint $q_{\max}$ must bound the
return: the learned returns span ${\approx}[19,48]$, so we set $q_{\max}{=}60$;
a naively small $q_{\max}$ sits below the distribution and silently handicaps
only the trapezoidal scheme.} surfaces here---on this task the return
distribution is smooth and well-scaled enough that the simple midpoint rule
already represents it adequately. We therefore adopt the midpoint rule (the
QR-DQN default) for all final models, as no alternative offers a significant
gain. Pilot runs with Simpson's rule and Beta-distorted grids reached the same
conclusion.

\begin{table}[htbp]
\caption{Quantile-positioning schemes, risk-neutral agent ($N{=}10$, $1000$ ep
$\times$ $3$ seeds; mean $\pm$ std on reward). The schemes are within cross-seed
noise on reward; midpoint is adopted as the simplest.}
\label{tab:qmode}
\begin{center}
\begin{tabular}{|l|c|c|c|}
\hline
\textbf{KPI} & \textbf{Midpoint} & \textbf{Gauss--Leg.} & \textbf{Trapez.} \\
\hline
Reward $\uparrow$                   & $1075.3\pm2.5$ & $1076.0\pm1.5$ & $1078.0\pm2.8$ \\
HOF/min $\downarrow$                & $1.89$  & $1.75$  & $1.60$  \\
RLF/min $\downarrow$                & $0.132$ & $0.130$ & $0.124$ \\
Capacity (bps/Hz) $\uparrow$        & $3.953$ & $3.955$ & $3.955$ \\
Reliability (\%) $\uparrow$         & $95.99$ & $95.94$ & $96.06$ \\
HO/min $\downarrow$                 & $15.31$ & $15.67$ & $15.40$ \\
PP/min $\downarrow$                 & $2.09$  & $2.09$  & $2.09$  \\
Prep/min $\downarrow$               & $1046$  & $1057$  & $1056$  \\
Res.\ reservation (\%) $\downarrow$ & $8.30$  & $8.39$  & $8.38$  \\
\hline
\end{tabular}
\end{center}
\end{table}

\subsection{Truncating the CVaR Head Is a Better Estimator}
\label{sec:results-trunc}
Where quantile positioning was inert, the truncated-CVaR head of
Section~\ref{sec:trunc-cvar} is decisive. We isolate it with two agents that
share the identical $\mathrm{CVaR}_{0.5}$ policy and differ \emph{only} in the
head: the \emph{truncated} head predicts $k=10$ quantiles over $[0,0.5]$, while
the \emph{full} head predicts all $N=20$ over $[0,1]$ and averages its lower ten
for selection ($1000$ ep $\times$ $3$ seeds, final backbone). Truncation
improves reward by $+8.3$ ($+8.8\sigma$, $p\approx0.001$; the per-seed ranges do
not overlap) and simultaneously lowers both failure rates (HOF $1.28$ vs.\
$1.84$, RLF $0.096$ vs.\ $0.138$) while raising capacity and reliability
(Table~\ref{tab:trunc}). The full head spends half its outputs on the
discarded upper lobe $[0.5,1]$; freeing that capacity for the lower tail that
actually drives the policy yields a \emph{better} CVaR estimate, not merely a
cheaper one. The truncated policy is in turn modestly more active---more handovers,
ping-pongs and preparations (HO $16.3$ vs.\ $15.6$, prep $1074$ vs.\ $1054$)---the
expected footprint of a more decisive risk-averse policy that pre-empts tail
failures, in the same direction as the RA-vs-RN contrast of
Section~\ref{sec:results-3way}. Finally, the smaller head trains ${\approx}20\%$
faster ($48$ vs.\ $61$\,min). Truncation is therefore both cheaper and a better
estimator, and we adopt it for the risk-aware agent.

\begin{table}[htbp]
\caption{Truncated vs.\ full CVaR head under the \emph{same} $\mathrm{CVaR}_{0.5}$
policy ($N{=}20$, $1000$ ep $\times$ $3$ seeds; mean $\pm$ std on reward, bold =
better). Truncating to the acted-upon lower tail is significantly better
\emph{and} cheaper; it trades modestly more mobility signalling for fewer
failures.}
\label{tab:trunc}
\begin{center}
\begin{tabular}{|l|c|c|}
\hline
\textbf{KPI} & \textbf{Truncated ($k{=}10$)} & \textbf{Full ($N{=}20$)} \\
\hline
Reward $\uparrow$                   & $\mathbf{1082.1\pm1.0}$ & $1073.7\pm1.3$ \\
HOF/min $\downarrow$                & $\mathbf{1.28}$  & $1.84$  \\
RLF/min $\downarrow$                & $\mathbf{0.096}$ & $0.138$ \\
Capacity (bps/Hz) $\uparrow$        & $\mathbf{3.976}$ & $3.949$ \\
Reliability (\%) $\uparrow$         & $\mathbf{96.24}$ & $95.98$ \\
HO/min $\downarrow$                 & $16.27$ & $\mathbf{15.61}$ \\
PP/min $\downarrow$                 & $2.09$  & $\mathbf{2.02}$  \\
Prep/min $\downarrow$               & $1074$  & $\mathbf{1054}$  \\
Res.\ reservation (\%) $\downarrow$ & $8.52$  & $\mathbf{8.37}$  \\
Train time (min) $\downarrow$       & $\mathbf{48}$ & $61$ \\
\hline
\end{tabular}
\end{center}
\end{table}

\subsection{Three-Way Comparison}
\label{sec:results-3way}
Table~\ref{tab:final} reports the learned agents at the canonical
$2000\times5$ budget against two references: our own hand-tuned LTM baseline,
run in the \emph{same} simulator and therefore directly comparable, and the
LTM contextual multi-armed bandit (CMAB) of~\cite{ainna}. The three headline
agents (DQN, RN, RA) form the matched-capacity comparison; the fourth QR-DQN
column is the single-quantile ($k{=}1$) risk-averse ablation of
Section~\ref{sec:results-bridge}. Two observations structure the result.

\textbf{Distributional learning helps, even with a risk-neutral policy.}
Moving from DQN to risk-neutral QR-DQN---same mean-based rule, same
architecture, only a distributional target---raises reward by $+14.1$
($11.7\sigma$) and cuts HOF by $36\%$ and RLF by $25\%$. This gain genuinely
needs \emph{multiple} quantiles: a risk-neutral QR-DQN restricted to a single
quantile is statistically indistinguishable from DQN, so the improvement is a
property of the distributional representation rather than the architecture
(Section~\ref{sec:results-bridge}).

\textbf{Risk-awareness adds a second, separable gain.} Switching the policy
from the mean to $\mathrm{CVaR}_{0.5}$ at identical capacity ($k=10$) adds a
further $+4.8$ reward ($5.1\sigma$) and cuts HOF by another $18\%$ and RLF by
another $13\%$. The two effects stack: the risk-aware agent improves reward,
HOF, RLF, capacity and reliability monotonically over both baselines.

\textbf{The trade-off is interpretable.} The same ordering that improves the
objective and safety KPIs \emph{increases} the activity/overhead KPIs: the
risk-aware agent performs more handovers ($15.5$ vs.\ DQN's $13.9$/min), more
cell preparations and more resource reservation. It is \emph{more willing to
act}, but its actions are far more reliable---HOF nearly halved relative to DQN
($-47\%$), RLF cut by $35\%$---and yield higher throughput. The agent buys
reliability and capacity with signalling activity (Fig.~\ref{fig:bars}).

\begin{table*}[t]
\caption{Final performance on the LTM HO task ($5$ seeds $\times$ $2000$ ep; mean $\pm$ std for the learned agents). \textbf{DQN}: scalar baseline. \textbf{RN}/\textbf{RA}: risk-neutral / risk-aware ($\mathrm{CVaR}_{0.5}$) QR-DQN at $k{=}10$. \textbf{RA-1q}: single-quantile risk-aware ablation. \textbf{LTM}: our hand-tuned LTM baseline, run in the \emph{same} simulator. \textbf{CMAB}: the LTM contextual multi-armed bandit of~\cite{ainna}. Arrows mark the desirable direction; bold marks the best \emph{learned} agent per row (baselines are shown for reference).}
\label{tab:final}
\begin{center}
\setlength{\tabcolsep}{5pt}
\begin{tabular}{|l|c|c|c|c|c|c|}
\hline
\textbf{KPI} & \textbf{DQN} & \textbf{RN} & \textbf{RA} & \textbf{RA-1q} & \textbf{LTM} & \textbf{CMAB}$^{*}$ \\
\hline
Reward $\uparrow$            & $1067.4\pm2.1$ & $1081.5\pm1.7$ & $\mathbf{1086.3\pm1.2}$ & $1080.6\pm1.2$ & $1055.5$ & --- \\
HOF/min $\downarrow$         & $2.233$ & $1.432$ & $\mathbf{1.176}$ & $1.469$ & $1.151$ & $0.97$ \\
RLF/min $\downarrow$         & $0.143$ & $0.107$ & $\mathbf{0.093}$ & $0.104$ & $0.065$ & $0.101$ \\
Capacity (bps/Hz) $\uparrow$ & $3.909$ & $3.951$ & $\mathbf{3.974}$ & $3.974$ & $3.753$ & $3.68$ \\
Reliability (\%) $\uparrow$  & $95.82$ & $96.10$ & $\mathbf{96.25}$ & $96.19$ & $95.25$ & $94.0$ \\
HO/min $\downarrow$          & $\mathbf{13.86}$ & $14.74$ & $15.47$ & $16.12$ & $11.34$ & $8.00$ \\
PP/min $\downarrow$          & $\mathbf{1.827}$ & $2.005$ & $1.946$ & $2.041$ & $3.625$ & $1.30$ \\
Prep/min $\downarrow$        & $\mathbf{1010.8}$ & $1038.8$ & $1057.6$ & $1072.9$ & $811.2$ & $20.0$ \\
Res.\ reservation (\%) $\downarrow$ & $\mathbf{8.022}$ & $8.244$ & $8.394$ & $8.515$ & $6.438$ & $6.5$ \\
\hline
\end{tabular}
\end{center}
{\footnotesize $^{*}$CMAB figures are reproduced from~\cite{ainna} and were obtained in that work's simulator. Our LTM baseline reproduces that work's published LTM KPIs to within $\approx 5\%$ on the radio and mobility metrics (resource reservation is the largest outlier, $\approx 13\%$), so the CMAB column is indicative rather than strictly comparable to our agents.}
\end{table*}

\begin{figure}[htbp]
\centerline{\includegraphics[width=\columnwidth]{master_bar_plots.png}}
\caption{The eight physical KPIs and the reward across the learned agents and
the LTM / LTM-CMAB baselines. Colour encodes the algorithm family (blue: the
LTM heuristic, shown for the paper's published numbers and our two simulators;
grey: LTM-CMAB; red: our learned agents), and a striped hatch marks values
reproduced from the paper.}
\label{fig:bars}
\end{figure}

\subsection{Where the Gain Comes From: a Quantile-Bridge Ablation}
\label{sec:results-bridge}
To localise the two gains we sweep the head from a single quantile up to the
full $k=10$ ($1000$ ep $\times$ $3$ seeds for the ablation; the single
risk-averse point is also re-run at the canonical $2000\times5$ budget).

\textbf{The distributional gain is a multi-quantile effect.} A risk-neutral
QR-DQN with a single quantile at $\tau=0.5$ reduces, by construction, to DQN:
the pinball loss at $\tau=0.5$ is a symmetric $\tfrac12$-Huber loss on the
median, collapsing the agent to scalar value learning. Empirically the two are
indistinguishable ($1063.8$ vs.\ $1060.9$, $0.8\sigma$). Reward then climbs
smoothly with the quantile count ($N{:}1\!\to\!10$ gives
$1063.8\!\to\!1075.4$), saturating by $N\approx5$.

\textbf{A single risk-averse quantile is not enough.} At the canonical budget a
single lower-tail quantile ($k=1$, $\tau=0.25$, $\alpha=0.5$) converges to
\emph{risk-neutral} performance (reward $1080.6$, HOF $1.469$)---tied with the
$k=10$ risk-neutral agent and $5.7$ reward ($7.8\sigma$) / $0.29$ HOF
($6.2\sigma$) below the full $k=10$ CVaR agent. Tellingly, it incurs the
\emph{highest} handover, preparation and reservation overhead of all agents
while delivering only risk-neutral outcomes: a single tail quantile is a
high-variance risk estimate that over-reacts without the payoff, whereas CVaR's
averaging over the lower $k=10$ quantiles is a stable estimate.

\textbf{Both axes are required.} Risk-aversion is necessary but not sufficient
on its own---it needs a stable, multi-quantile (CVaR) estimate. Neither a
single risk-averse quantile nor a risk-neutral $10$-quantile model reaches the
champion; only their combination does.

\subsection{Comparison with the LTM Baseline and Contextual Bandits}
\label{sec:results-baselines}
\textbf{Versus our LTM baseline (same simulator).} Against the hand-tuned LTM
heuristic run in our own simulator---an apples-to-apples reference free of the
cross-simulator caveat below---the risk-aware agent improves every user-facing
KPI: $+5.9\%$ capacity ($3.974$ vs.\ $3.753$\,bps/Hz), $+1.0$ point of
reliability ($96.25\%$ vs.\ $95.25\%$) and a $46\%$ lower ping-pong rate
($1.95$ vs.\ $3.63$/min), while matching the heuristic's handover-failure rate
($1.176$ vs.\ $1.151$); its reward is $+30.9$ points higher ($1086.3$ vs.\
$1055.5$). As in the three-way comparison it pays for this with more
control-plane activity---more handovers ($15.5$ vs.\ $11.3$/min), preparations
and reservations---and a higher RLF ($0.093$ vs.\ $0.065$): the policy is more
willing to act, but its actions are more reliable and higher-throughput.

\textbf{Versus the contextual bandit (CMAB).} The CMAB of~\cite{ainna} adds a
target pre-selection layer on top of LTM; we reproduce its numbers from that
work. Because they were measured in a different simulator
(Table~\ref{tab:final}, footnote), we read them as indicative. The risk-aware
agent is clearly ahead on the user-facing quality KPIs---$+8.0\%$ capacity
($3.974$ vs.\ $3.68$), $+2.25$ points of reliability ($96.25\%$ vs.\ $94.0\%$)
and a lower RLF ($0.093$ vs.\ $0.101$). CMAB, in turn, is far more frugal with
signalling: by pre-selecting targets it issues dramatically fewer preparations
($20$ vs.\ $1058$/min) and fewer handovers ($8.0$ vs.\ $15.5$/min) and
ping-pongs. The two mechanisms are therefore complementary: CMAB, like the
orthogonal RIS and LMMSE-prediction enhancements of~\cite{ainna}, acts on
target pre-selection and signal quality rather than on the value estimate, so
it could be stacked on top of the distributional-RL policy; we leave that
combination to future work.

% =====================================================================
\section{Discussion, Limitations, and Future Work}
\label{sec:discussion}
\textbf{Limitations.} Several caveats bound our claims. (i)~Training and
evaluation use the \emph{same} fixed $1{,}000$-trajectory set with no held-out
split; this is the appropriate protocol for the question we ask---how well a
policy exploits a fixed propagation environment, matching the evaluation
of~\cite{ainna}---but it does not measure generalisation to unseen geometry.
(ii)~We study a single deployment and channel-gain dataset. (iii)~The
reward-shaping coefficients are inherited from~\cite{ainna} and not re-optimised
here; $\kappa$, the network width and the quantile placement, by contrast, are
all swept on the final backbone (Sections~\ref{sec:setup},~\ref{sec:results-qmode})
and leave the operating point within noise. (iv)~CVaR is the only risk measure we
consider; other coherent or distortion risk measures may trade the KPIs off
differently.

\textbf{Future work.} Natural extensions include implicit and
fully-parameterised quantile networks (IQN~\cite{iqn}, FQF~\cite{fqf}), which
learn the quantile fractions themselves and would subsume the
quantile-positioning study; constrained or multi-objective formulations that
bound HOF or RLF directly rather than through a scalar reward; and validation
on higher-fidelity or over-the-air deployments. Because the RIS,
LMMSE-prediction and contextual-bandit enhancements of~\cite{ainna} act on
signal quality and target pre-selection rather than on the value estimate, they
are complementary to our risk-aware policy and could be stacked on top of it.
Finally, it remains to confirm that the distributional and quantile-positioning
findings transfer to standard finite-action RL benchmarks such as the
Atari~2600 suite.

% =====================================================================
\section{Conclusion}
\label{sec:conclusion}
On the LTM HO task, distributional RL delivers two cleanly separable
improvements over scalar DQN: a large gain from the distributional learning
target itself (DQN $\to$ RN), and a second, smaller-but-significant gain from a
CVaR-based risk-aware policy (RN $\to$ RA) whose optimal level
($\alpha\approx0.5$) we establish with a de-confounded sweep. The resulting
risk-aware QR-DQN attains the best reward and the best safety and throughput
profile of any learned agent, halving the handover-failure rate of DQN and
exceeding our LTM baseline on capacity, reliability and ping-pong rate.
A quantile-bridge ablation confirms both ingredients are necessary: a single
risk-neutral quantile collapses to DQN, and even a single risk-averse quantile
only reaches the risk-neutral $10$-quantile level---the combined gain requires
the full $\sim\!10$-quantile CVaR representation.

% =====================================================================
\section*{Acknowledgment}
\PLACEHOLDER{Funding and acknowledgements.}

% =====================================================================
\begin{thebibliography}{00}
\bibitem{ainna} A. Y. Moreno-Locubiche \emph{et al.}, ``Contextual bandits and reconfigurable intelligent surfaces for predictive LTM handover decisions,'' \PLACEHOLDER{venue, year, pages}.
\bibitem{c51} M. G. Bellemare, W. Dabney, and R. Munos, ``A distributional perspective on reinforcement learning,'' in \emph{Proc. ICML}, 2017.
\bibitem{qrdqn} W. Dabney, M. Rowland, M. G. Bellemare, and R. Munos, ``Distributional reinforcement learning with quantile regression,'' in \emph{Proc. AAAI}, 2018.
\bibitem{iqn} W. Dabney, G. Ostrovski, D. Silver, and R. Munos, ``Implicit quantile networks for distributional reinforcement learning,'' in \emph{Proc. ICML}, 2018.
\bibitem{fqf} D. Yang, L. Zhao, Z. Lin, T. Qin, J. Bian, and T.-Y. Liu, ``Fully parameterized quantile function for distributional reinforcement learning,'' in \emph{Proc. NeurIPS}, 2019.
\bibitem{rockafellar} R. T. Rockafellar and S. Uryasev, ``Optimization of conditional value-at-risk,'' \emph{J. Risk}, vol.~2, pp.~21--41, 2000.
\bibitem{chow2014} Y. Chow and M. Ghavamzadeh, ``Algorithms for CVaR optimization in MDPs,'' in \emph{Proc. NeurIPS}, 2014.
\bibitem{tamar2015} A. Tamar, Y. Glassner, and S. Mannor, ``Optimizing the CVaR via sampling,'' in \emph{Proc. AAAI}, 2015.
\bibitem{dqn} V. Mnih \emph{et al.}, ``Human-level control through deep reinforcement learning,'' \emph{Nature}, vol.~518, pp.~529--533, 2015.
\end{thebibliography}

\vspace{6pt}
\textcolor{red}{\textbf{Note:} red \texttt{[Placeholder]}/\texttt{[TODO]} markers and this note must be removed before submission. Figures referenced by \texttt{\textbackslash includegraphics} must be uploaded alongside this file.}

\end{document}
